% Part: incompleteness
% Chapter: representability-in-q
% Section: introduction

\documentclass[../../../include/open-logic-section]{subfiles}

\begin{document}

\olfileid{inc}{req}{int}
\olsection{ভূমিকা}

যেসব তত্ত্বে গণনসাধ্য অপেক্ষক সম্বন্ধে মৌলিক তথ্য প্রকাশ ও প্রমাণ করা যায়,
অসম্পূর্ণতা উপপাদ্যগুলি সেই তত্ত্বগুলিতে প্রযোজ্য। এমনই একটি অত্যন্ত ন্যূনতম
তত্ত্বের বর্ণনা দেব; তার নাম ``$\Th{Q}$'' (অথবা কখনও রাফায়েল রবিনসনের নাম
অনুসারে ``রবিনসনের $Q$'')। কোনো অপেক্ষক $\Th{Q}$-তে
\emph{উপস্থাপনযোগ্য} হওয়ার অর্থ বলব, এবং তারপর নিচের কথাটি প্রমাণ করব:
\begin{quote}
  কোনো অপেক্ষক $\Th{Q}$-তে উপস্থাপনযোগ্য যদি এবং কেবল যদি সেটি গণনসাধ্য।
\end{quote}
এক দিক থেকে এটি আমাদের গণনসাধ্যতার আরেকটি মডেল দেয়। কিন্তু থামার সমস্যাকে
সেটিতে হ্রাস করে আমরা এ-ও দেখাব যে
$\Setabs{!A}{\Th{Q} \Proves !A}$ সেটটি নির্ণেয় নয়। কাজ শেষ হতে হতে এর
চেয়েও অনেক শক্তিশালী ফল প্রমাণিত হবে।

$\Th{Q}$-এর ভাষা পাটিগণিতের ভাষা; !!{identity}-সহ প্রথম-ক্রমের যুক্তিবিদ্যার অন্য
স্বতঃসিদ্ধ ও বিধির সঙ্গে ব্যবহার করার জন্য $\Th{Q}$ নিচের স্বতঃসিদ্ধগুলি নিয়ে
গঠিত:
\begin{align*}
& \lforall[x][\lforall[y][(\eq[x'][y'] \lif \eq[x][y])]] \tag{$!Q_1$}\\
& \lforall[x][\eq/[\Obj 0][x']] \tag{$!Q_2$}\\
& \lforall[x][(\eq[x][\Obj 0] \lor \lexists[y][\eq[x][y']])] \tag{$!Q_3$}\\
& \lforall[x][\eq[(x + \Obj 0)][x]] \tag{$!Q_4$}\\
& \lforall[x][\lforall[y][\eq[(x + y')][(x + y)']]] \tag{$!Q_5$}\\
& \lforall[x][\eq[(x \times \Obj 0)][\Obj 0]] \tag{$!Q_6$}\\
& \lforall[x][\lforall[y][\eq[(x \times y')][((x \times y) + x)]]] \tag{$!Q_7$}\\
& \lforall[x][\lforall[y][(x < y \liff \lexists[z][\eq[(z' + x)][y]])]] \tag{$!Q_8$}
\end{align*}
প্রত্যেক স্বাভাবিক সংখ্যা $n$-এর জন্য সংখ্যাপদ $\num{n}$-কে
$\Obj{0}^{\prime\prime\ldots\prime}$ পদ হিসেবে সংজ্ঞায়িত করি, যেখানে মোট
$n$-টি টিক-চিহ্ন আছে। তাই $\num{0}$ হলো একা !!{constant}~$\Obj{0}$,
$\num{1}$ হলো $\Obj{0}'$, $\num{2}$ হলো $\Obj{0}''$, ইত্যাদি।

পাটিগণিতের তত্ত্ব হিসেবে $\Th{Q}$ \emph{অত্যন্ত} দুর্বল; যেমন,
$\lforall[x][\eq/[x][x']]$ অথবা $\lforall[x][\lforall[y][(x + y) = (y +
    x)]]$-এর মতো অতি সরল তথ্যও এতে প্রমাণ করা যায় না। তবে আমরা দেখব,
$\Th{Q}$ এত আকর্ষণীয় হওয়ার বড় কারণই হলো এর এই দুর্বলতা। বস্তুত অসম্পূর্ণতা
উপপাদ্য প্রযোজ্য হওয়ার জন্য এটি ঠিক যতটুকু শক্তিশালী হওয়া দরকার, প্রায়
ততটুকুই। $\Th{Q}$ আকর্ষণীয় হওয়ার আরেকটি কারণ হলো এর স্বতঃসিদ্ধের সেট
\emph{সসীম}।

$\Th{Q}$-এর সঙ্গে আরোহের একটি ছক যোগ করলে $\Th{Q}$-এর চেয়ে শক্তিশালী তত্ত্ব
\emph{পেয়ানো পাটিগণিত} $\Th{PA}$ পাওয়া যায়:
\[
(!A(\Obj 0) \land \lforall[x][(!A(x) \lif !A(x'))]) \lif \lforall[x][!A(x)]
\]
এখানে $!A(x)$ যেকোনো সূত্র। $!A(x)$-এ $x$ ছাড়া অন্য কোনো মুক্ত
!!{variable}s থাকলে, সবগুলিকে বদ্ধ করতে সামনে সর্বজনীন পরিমাণসূচক যোগ করি
(যাতে আরোহ-ছকের সংশ্লিষ্ট নিদর্শনটি !!a{sentence} হয়)। যেমন, $!A(x, y)$-এ
!!{variable}~$y$-ও মুক্ত থাকলে সংশ্লিষ্ট নিদর্শনটি হলো
\[
\lforall[y][((!A(\Obj 0) \land \lforall[x][(!A(x) \lif !A(x'))]) \lif
  \lforall[x][!A(x)])]
\]
আরোহ-ছকের নিদর্শনগুলি ব্যবহার করে $\Th{Q}$-এর স্বতঃসিদ্ধগুলি থেকে যতটা প্রমাণ
করা যায়, $\Th{PA}$-এর স্বতঃসিদ্ধগুলি থেকে তার চেয়ে অনেক বেশি প্রমাণ করা
যায়। বস্তুত, স্বাভাবিক সংখ্যা সম্বন্ধে পেয়ানো পাটিগণিতে প্রমাণ করা যায় না
এমন ``স্বাভাবিক'' বক্তব্য খুঁজতে বেশ পরিশ্রম করতে হয়!{}

\begin{defn}
\ollabel{defn:representable-fn}
  স্বাভাবিক সংখ্যা থেকে স্বাভাবিক সংখ্যায় অপেক্ষক
  $f(x_0,\ldots,x_k)$-কে {$\Th{Q}$-তে \emph{উপস্থাপনযোগ্য}} বলা হয়, যদি
  এমন একটি সূত্র $!A_f(x_0,\dots,x_k,y)$ থাকে যাতে যখনই
  $f(n_0,\dots,n_k) = m$, তখন $\Th{Q}$ নিচের দুটিকে প্রমাণ করে:
\begin{enumerate}
\item\ollabel{defn:rep:a} $!A_f(\num{n_0}, \dots, \num{n_k}, \num{m})$
\item\ollabel{defn:rep:b} $\lforall[y][(!A_f(\num{n_0}, \dots,
\num{n_k}, y) \lif \num{m} = y)]$।
\end{enumerate}
\end{defn}

সংজ্ঞাটি অন্যভাবেও বলা যায়; যেমন, সমতুল্যভাবে আমরা চাইতে পারি যে $\Th{Q}$
$\lforall[y][(!A_f(\num{n_0}, \dots, \num{n_k}, y) \liff
\eq[y][\num{m}])]$ প্রমাণ করুক।

\begin{thm}
\ollabel{thm:representable-iff-comp}
কোনো অপেক্ষক $\Th{Q}$-তে উপস্থাপনযোগ্য যদি এবং কেবল যদি সেটি গণনসাধ্য।
\end{thm}

উপপাদ্যটি প্রমাণের দুটি দিক আছে। বাক্যগঠনের পাটিগণিতায়ন হয়ে গেলে বাম থেকে
ডান দিকটি বেশ সরল। অন্য দিকে কিছু বেশি কাজ দরকার। মূল ধারণাটি এই: আমরা
``গণনসাধ্য''-কে সুনির্দিষ্ট করার উপায় হিসেবে ``সাধারণ পুনরাবৃত্ত'' বেছে নিই
এবং দেখাই যে প্রত্যেক সাধারণ পুনরাবৃত্ত অপেক্ষক $\Th{Q}$-তে উপস্থাপনযোগ্য।
মনে করো, কোনো অপেক্ষক তখনই সাধারণ পুনরাবৃত্ত, যখন শূন্য অপেক্ষক $\Zero$,
উত্তরসূরি অপেক্ষক~$\Succ$ এবং অভিক্ষেপ অপেক্ষক~$\Proj{n}{i}$ থেকে মিশ্রণ,
আদিম পুনরাবৃত্তি ও নিয়মিত ন্যূনীকরণ ব্যবহার করে তাকে সংজ্ঞায়িত করা যায়।
তাই প্রত্যেক সাধারণ পুনরাবৃত্ত অপেক্ষক $\Th{Q}$-তে উপস্থাপনযোগ্য দেখানোর
একটি উপায় হলো: মৌলিক অপেক্ষকগুলি উপস্থাপনযোগ্য দেখানো, এবং কয়েকটি অপেক্ষক
উপস্থাপনযোগ্য হলে তাদের থেকে মিশ্রণ, আদিম পুনরাবৃত্তি ও নিয়মিত ন্যূনীকরণে
সংজ্ঞায়িত অপেক্ষকগুলিও উপস্থাপনযোগ্য দেখানো। অন্যভাবে বললে, মৌলিক
অপেক্ষকগুলি উপস্থাপনযোগ্য এবং উপস্থাপনযোগ্য অপেক্ষকগুলি মিশ্রণ, আদিম
পুনরাবৃত্তি ও নিয়মিত ন্যূনীকরণের ``অধীনে বদ্ধ''—এই দুটি কথা দেখানো যায়।
তাতে প্রত্যেক সাধারণ পুনরাবৃত্ত অপেক্ষকের উপস্থাপনযোগ্যতা নিশ্চিত হয়।

কিন্তু উপস্থাপনযোগ্য অপেক্ষকগুলি আদিম পুনরাবৃত্তির অধীনে বদ্ধ দেখানোর ধাপটি
কঠিন। সেই ধাপ এড়াতে আমরা প্রথমে দেখাই যে আদিম পুনরাবৃত্তি ছাড়াই আসলে কাজ
চলে। অর্থাৎ, শুধু মিশ্রণ ও নিয়মিত ন্যূনীকরণ ব্যবহার করে মৌলিক অপেক্ষক থেকে
প্রত্যেক সাধারণ পুনরাবৃত্ত অপেক্ষক সংজ্ঞায়িত করা যায়। এর জন্য দেখাই, একটি
নির্দিষ্ট নিয়মিত ন্যূনীকরণ দিয়েই আদিম পুনরাবৃত্তি করা যায়। তবে তা কাজে
লাগাতে আরও কয়েকটি মৌলিক অপেক্ষক যোগ করতে হবে: যোগ, গুণ এবং পরিচয়-সম্বন্ধের
চরিত্রসূচক অপেক্ষক~$\Char{=}$। এরপর \emph{এই} সব মৌলিক অপেক্ষক $\Th{Q}$-তে
উপস্থাপনযোগ্য এবং উপস্থাপনযোগ্য অপেক্ষকগুলি মিশ্রণ ও নিয়মিত ন্যূনীকরণের
অধীনে বদ্ধ—এই দুটি কথা দেখিয়ে উপপাদ্যটি প্রমাণ করা যাবে।

\end{document}
